\documentclass[a4paper, top=10mm]{article}
%for writing from the top
\usepackage{fullpage}
%for math
\usepackage{amsmath}
\usepackage{mathrsfs}
\usepackage{amsthm}
%for images
\usepackage{graphicx}
%for color
\usepackage{xcolor}
%for title
\title{\textbf{\huge{Scores de rugby}}}
\author{Enigma n\textsuperscript{o}4}
\date{16\textsuperscript{th} June 2026}

\newtheorem*{hint}{Hint}

\addtolength{\voffset}{-2cm}
\addtolength{\textheight}{5cm}


\begin{document}
	\maketitle
	
	\Large
	Au rugby, 10 est le plus petit score pouvant etre obtenu de deux facons differentes (deux essais non transformes, car $5+5=10$; ou un essai transforme et un drop, car $5+2+3=10$).
	
	Quel est le plus petit score pouvant etre obtenu de 3 facons differentes?
	
	\vspace{3cm}
	
	In rugby, 10 is the lowest score that can be obtained in two different ways (two unconverted tries, because $5 + 5 = 10$; or one converted try and one drop goal, because $5 + 2 + 3 = 10$).
	
	What is the lowest score that can be obtained in three different ways?
	
	
	% 15 = 5+5+5, 3+3+3+3+3, 7+5+3
	
\end{document}